\subsection{Local Rings}

PROPOSITION 1. --- Let A be a nonzero ring, and let r be the set of noninvertible elements of A. The following properties are equivalent:

\begin{enumerate}
\def\labelenumi{(\roman{enumi})}
\item
  The set \(\mathfrak{r}\) is a two-sided ideal of A.
\item
  The set \(\mathfrak{r}\) is stable under addition.
\item
  The ring A has a unique maximal left ideal.
\item
  For every \(a \in \mathbf{A}\) , one of the elements \(a\) and \(1 - a\) is invertible.
\item
  For every \(a \in \mathbf{A}\) , one of the elements \(a\) and \(1 - a\) is left invertible.
\end{enumerate}

The implication (i) \(\Rightarrow\) (ii) follows from the definition of an ideal. Since 1 does not belong to \(\mathfrak{r}\) , we have (ii) \(\Rightarrow\) (iv).

We have \(\mathfrak{r} \neq \mathrm{A}\) , and the set \(\mathfrak{r}\) contains every left ideal of \(\mathrm{A}\) not equal to \(\mathrm{A}\) . If \(\mathfrak{r}\) is a left ideal of \(\mathrm{A}\) , it is therefore the unique maximal left ideal of \(\mathrm{A}\) . This proves that (i) implies (iii).

Suppose that A has a unique maximal left ideal m. Take \(b \in A - m\) . The left ideal Ab is not contained in any maximal left ideal of A, hence is equal to A (I, §8, No.~6, p.~104, Theorem 1), and b is left invertible. For every \(a \in A\) , one of the elements a and 1 - a belongs to A - m because m is an ideal that does not contain 1. Thus, (iii) implies (v).

Suppose that property (v) holds. Let b be a left invertible element of A. Let \(c \in A\) be such that cb = 1. We have \((1 - bc)b = 0\) and \(b \neq 0\) ; hence, 1 - bc is not left invertible. By property (v), bc is left invertible and, a fortiori, c is left invertible. But then c is invertible; b is its inverse, so b is invertible. It follows that (v) implies (iv).

It remains to prove that (iv) implies (i). Suppose that (iv) holds. Then \(\mathfrak{r}\) is a two-sided ideal of A by the following assertions a) through d):

\begin{enumerate}
\def\labelenumi{\alph{enumi})}
\item
  We have \(0 \in \mathfrak{r}\) because the ring A is nonzero.
\item
  The product of two elements of A, one belonging to r and the other to A-r, belongs to r.
\item
  The set \(\mathfrak{r}\) is stable under addition.
\end{enumerate}

Indeed, let \(a\) and \(b\) be elements of \(\mathfrak{r}\) such that \(s = a + b\) is invertible. By b), the elements \(s^{-1}a\) and \(s^{-1}b\) of A belong to \(\mathfrak{r}\) ; since \(s^{-1}b = 1 - s^{-1}a\) , this contradicts the assumption that (iv) holds.

\begin{enumerate}
\def\labelenumi{\alph{enumi})}
\setcounter{enumi}{3}
\tightlist
\item
  The set \(\mathfrak{r}\) is stable under multiplication.
\end{enumerate}

Indeed, let \(a\) and \(b\) be two elements of \(\mathfrak{r}\) . The element \(a' = -a(1 - b)\) belongs to \(\mathfrak{r}\) by \(\mathrm{b}\) ), so that the element \(ab\) , which is equal to \(a + a'\) , belongs to \(\mathfrak{r}\) by \(\mathrm{c}\) ; assertion d) follows.

DEFINITION 1. --- A local ring is a nonzero ring that has the equivalent properties of Proposition 1.

A ring A is local if and only if the opposite ring \(A^{o}\) is local.

If A is a local ring, then the set r of noninvertible elements of A is a two-sided ideal of A; it contains every left or right ideal of A not equal to A. The ring \(A/r\) is therefore a field, which we call the residue field of A. The set r is the unique maximal left (resp. right, two-sided) ideal of A; we simply say that r is the maximal ideal of A.

Examples. --- 1) Every field is a local ring.

\begin{enumerate}
\def\labelenumi{\arabic{enumi})}
\setcounter{enumi}{1}
\tightlist
\item
  Let A be a nonzero ring in which every element is invertible or nilpotent. Then A is a local ring. Indeed, if \(a \in A\) is not invertible, then by assumption, there exists an integer \(n \geqslant 0\) such that \(a^{n+1} = 0\) , and 1 - a has inverse \(1 + a + \cdots + a^{n}\) .
\end{enumerate}

*3) Let \(\mathrm{X}\) be a \(\mathrm{C}^r\) manifold (VAR, R, 5.1.5) and \(x\) a point of \(\mathrm{X}\) . Let \(\mathcal{O}_x\) be the ring of germs at \(x\) of \(\mathrm{C}^r\) functions with values in the field of scalars \(\mathrm{K}\) . Then \(\mathcal{O}_x\) is a commutative local ring, and its maximal ideal consists of the germs of the functions that are zero at \(x\) .

\begin{enumerate}
\def\labelenumi{\arabic{enumi})}
\setcounter{enumi}{3}
\item
  Let A be a commutative local ring and \(B = A[[X_{i}]]_{i \in I}\) an algebra of formal power series with coefficients in A (III, §2, No.~11, p.~456). By Proposition 6 of IV, §4, No.~4, p.~30, the ring B is local, and its maximal ideal consists of the formal power series with constant term in the maximal ideal of A. In particular, if A is a field, the maximal ideal of \(A[[X_{i}]]_{i \in I}\) consists of the formal power series with constant term zero.
\item
  Let p be a prime number. We denote by \(\mathbf{Z}_{(p)}\) the subring of the field Q of rational numbers consisting of the fractions a/b with \(a \in Z\) , \(b \in Z\) , and b not divisible by \(p\) *(cf.~Comm. Alg., II, §2, No.~1, p.~60)*. Then \(\mathbf{Z}_{(p)}\) is a commutative local ring, with maximal ideal \(p\mathbf{Z}_{(p)}\) . The ring \(Z_{p}\) of p-adic integers (V, §12, No.~3, p.~96) is a commutative local ring, with maximal ideal \(pZ_{p}\) (VIII, p.~40, Exercise 9).
\item
  Let K be a commutative field of characteristic p \textgreater{} 0 and G a p-group (I, §6, No.~5, p.~76, Definition 9). The algebra K{[}G{]} of the group G over K (III, §2, No.~6, p.~446) is a local ring; its maximal ideal is the set of elements \((a_{g})_{g \in \mathrm{G}}\) of K{[}G{]} such that \(\sum_{g \in G} a_{g} = 0\) (VIII, p.~41, Exercise 10).
\end{enumerate}

